%% kapa: exam=1 round=21 year=2010 period=1 form=A subject=civil_law subject_ko=민법 seq=28 number=28
%% source: Q-Net 감정평가사 기출문제 (cst00309) artlSeq=5024181 file=extracted/1008/5024181/2010년 제21회 감정평가사 1차시험 기출문제/1교시.hwp
%% conversion: kordoc 4.24.51 -> tools/gen1.py (auto); review=proofread-sonnet; flags=none
\begin{question}{28}
지상권의 지료 지급에 관한 설명으로 옳지 않은 것은? (다툼이 있으면 판례에 의함)
\begin{choices}
  \item 법정지상권의 지료가 정해졌다는 증명이 없다면, 토지소유자는 지상권자가 2년 이상의 지료를 지급하지 않았음을 이유로 지상권의 소멸을 청구할 수 없다.
  \item 당사자간에 지료의 지급을 서면으로 약정하였다면 그 등기 없이도 제3자에게 대항할 수 있다.
  \item 지상권자의 지료지급 연체가 토지소유권의 양도 전후에 걸쳐 이루어진 경우, 토지양수인에 대한 연체기간이 2년이 되지 않는다면 양수인은 지상권소멸청구를 할 수 없다.
  \item 법정지상권이 성립한 경우, 그 지료액에 대하여 당사자간에 협의가 되지 않으면 당사자의 청구에 의해 법원이 이를 정한다.
  \item 분묘기지권을 시효취득한 경우에는 임료상당의 지료를 지급할 필요가 없는 것이 원칙이다.
\end{choices}
\end{question}
\answer{2}
